\documentclass[12pt]{article}
\usepackage{geometry}
\geometry{letterpaper, margin=1in}
\usepackage{hyperref}
\usepackage[usenames,dvipsnames]{xcolor}
\usepackage{microtype}
\usepackage{enumitem}

\definecolor{accent}{RGB}{26, 60, 110}

\begin{document}

%----------HEADER----------
\begin{center}
  {\LARGE\textbf{\textcolor{accent}{SITONG LI}}} \\[4pt]
  \small Ann Arbor, MI
  \;\textbullet\;
  \href{mailto:ruc.sitongli@gmail.com}{ruc.sitongli@gmail.com}
  \;\textbullet\;
  (952) 201-4026
\end{center}

\vspace{1em}

\noindent May 23, 2026

\vspace{1em}

\noindent Dear Bristol Myers Squibb AI Venture Studio Hiring Team,

\vspace{1em}

I am writing to apply for the AI Engineer position within Bristol Myers Squibb's AI Venture Studio delivery team. I am currently a full-time AI engineer at the Ross School of Business, University of Michigan, where I design and maintain the data pipelines, LLM integrations, and agentic AI systems that power large-scale research projects for cross-institutional faculty teams. The Venture Studio's model --- rapid 12-week project cycles, matrixed stakeholders, and a mandate to build reliable data foundations that let AI systems move faster --- is precisely the environment I thrive in, and I am eager to apply that experience to work that can directly improve patients' lives.

At Michigan Ross, I have built two production-grade systems end to end. The first is a scalable entity resolution pipeline that integrates LLMs, fuzzy string matching, and word embeddings to link records across heterogeneous administrative databases --- a project requiring careful attention to data quality validation at system boundaries and iterative calibration against ground-truth labels. The second is an agentic computer vision pipeline that extracts and classifies demographic attributes from applicant images at scale, with automated reporting to support downstream stakeholder analysis. Both systems were built to serve researchers who needed fast, reliable outputs within tight project timelines, which has directly trained the engineering discipline and stakeholder communication skills that drive the Venture Studio's 12-week delivery cycles.

My prior research engineering work covers every area of the job description. At Washington University in St.\ Louis, I built LLM-powered extraction workflows using GPT-4 to classify topics and extract structured fiscal policy indicators from 100,000+ lines of congressional testimony, and used the multimodal LLM Donut to construct a reproducible data ingestion pipeline from scanned archival documents. Ensuring extraction quality required prompt engineering, structured output validation, and iterative evaluation --- directly analogous to the LLM observability and evaluation pipelines BMS builds using tools like LangSmith. At the University of Minnesota, I built an agentic GPT-4 search bot integrated with the Google Search API to automate structured data enrichment across 1,000+ regulated entities --- a multi-step, tool-integrated agent workflow engineered from scratch in Python. In the same group, I designed an OCR-based ETL pipeline to convert ten years of regulatory correspondence into structured, analysis-ready data, then matched the resulting dataset with EPA enforcement records. In a separate project, I built a web crawler to ingest 10,000+ SEC filings, applied BERT classifiers across 500,000+ sentences, and implemented data quality checks throughout the pipeline to ensure reliable downstream outputs.

Across all of these projects, my work centers on exactly what the Venture Studio does: building robust data foundations and agentic AI systems that enable rapid, reliable delivery. I work in Python with standard data libraries, manage infrastructure on AWS, use Git and CI practices as a matter of course, and bring direct experience with LLM integration via the OpenAI and Anthropic APIs. I prototype quickly and engineer for production. I am a direct, collaborative communicator who has worked with matrixed stakeholders across institutions and consistently delivered within tight timelines. I am drawn to BMS because the Venture Studio's combination of cutting-edge AI engineering and measurable patient impact is the kind of work I want to build a career around.

I would welcome the opportunity to discuss how my background fits the Venture Studio's current project portfolio.

\vspace{1em}

\noindent Sincerely,\\[0.5em]
\textbf{Sitong Li}

\end{document}
